% GENERATED FILE. Edit canonical lessons, then run npm run generate:books.
% canonical-source-hash: fnv1a64:2720265718580c33
% canonical-lessons: ZH-C259-zhiye, ZH-C259-zuojia, ZH-C259-lushi, ZH-C259-jizhe, ZH-C259-diangong, ZH-R259-review-254-259

\chapter{What Do You Do?}
\label{ch:zh-what-do-you-do}
\hlchaptermodality{259}

\begin{chapteropening}
\textbf{By the end of this chapter:} \emph{I can ask what someone does for a living and name an author, a lawyer, a journalist and an electrician.}

Complete the last lesson of chapter 259: I can ask what someone does for a living and name an author, a lawyer, a journalist and an electrician.
\end{chapteropening}

\section[zhíyè]{zhíyè — occupation; profession}

\label{lesson:ZH-C259-zhiye}



\begin{quote}
Before the new one: say the Mandarin for witty, then the Mandarin for personality.
\end{quote}

\subsection*{You'll want to know: zhíyè}
\emph{zhíyè} — \textquotedblleft{}occupation; profession\textquotedblright{}.

\begin{quote}
\emph{nǐ de zhíyè shì shénme} — What do you do for a living?
\end{quote}

Its characters come later; for now it is a word for the ear.

\begin{cousinweb}
\emph{zhíyè} is built from \emph{zhí} 'a duty' and \emph{yè} 'a task', both inherited from Middle Chinese, \emph{tsyik} and \emph{ngjaep} (Guangyun, in Baxter's transcription), which quotes two older dictionaries for the first: 'to be in charge' and 'work'.
\end{cousinweb}

\begin{grammarlens}[title={What you've built}]
One more word for saying what I do.
\end{grammarlens}

\subsection*{Your turn}
\begin{itemize}
\raggedright
  \item \emph{Say it:} \emph{zhíyè}
  \item \emph{Say it:} \emph{zhíyè}, once more
  \item \emph{Say it:} \emph{nǐ de zhíyè shì shénme} — What do you do for a living?
  \item \emph{Recall:} say the Mandarin for witty, then the Mandarin for personality, then say \emph{zhíyè} again
\end{itemize}

\begin{tcolorbox}[breakable,colback=teal!4,colframe=teal!35!black,title={Before you move on}]
What does zhíyè mean? (\textquotedblleft{}Occupation; profession\textquotedblright{}.) Say it once more.
\end{tcolorbox}
\section[\texorpdfstring{\zh{作家} (zuòjiā)}{作家 (zuòjiā)}]{\zh{作家} (zuòjiā) — writer; author}

\label{lesson:ZH-C259-zuojia}



\begin{quote}
Before the new one: say the Mandarin for personality, then the Mandarin for an occupation.
\end{quote}

\subsection*{You'll want to know: \zh{作家}}
\textbf{\zh{作家}} — \emph{zuòjiā} — \textquotedblleft{}writer; author\textquotedblright{}.

\begin{quote}
\textbf{\zh{他是作家写了十本书}} \emph{tā shì zuòjiā xiěle shí běn shū} — He's an author and has written ten books.
\end{quote}

\begin{cousinweb}
\textbf{\zh{作家}} (\emph{zuòjiā}) is built from \zh{作} 'to make' and \zh{家}, both inherited from Middle Chinese, \emph{tsak} and \emph{kae} (Guangyun, in Baxter's transcription), which glosses the second 'a dwelling'. \zh{家} also names a master of a craft, a sense the dictionary does not give, so how it came is not settled.
\end{cousinweb}

\begin{grammarlens}[title={What you've built}]
One more word for saying what I do.
\end{grammarlens}

\subsection*{Your turn}
\begin{itemize}
\raggedright
  \item \emph{Say it:} \emph{zuòjiā}
  \item \emph{Say it:} \emph{zuòjiā}, once more
  \item \emph{Say it:} \emph{tā shì zuòjiā xiěle shí běn shū} — He's an author and has written ten books.
  \item \emph{Recall:} say the Mandarin for personality, then the Mandarin for an occupation, then say \emph{zuòjiā} again
\end{itemize}

\begin{tcolorbox}[breakable,colback=teal!4,colframe=teal!35!black,title={Before you move on}]
What does \zh{作家} mean? (\textquotedblleft{}Writer; author\textquotedblright{}.) Say it once more.
\end{tcolorbox}
\section[lǜshī]{lǜshī — lawyer}

\label{lesson:ZH-C259-lushi}



\begin{quote}
Before the new one: say the Mandarin for an occupation, then the Mandarin for an author.
\end{quote}

\subsection*{You'll want to know: lǜshī}
\emph{lǜshī} — \textquotedblleft{}lawyer\textquotedblright{}.

\begin{quote}
\emph{wǒ jiějie shì lǜshī} — My sister is a lawyer.
\end{quote}

Its characters come later; for now it is a word for the ear.

\begin{cousinweb}
\emph{lǜshī} is built from \emph{lǜ} 'the law' and \emph{shī} 'a master', both inherited from Middle Chinese, \emph{lwit} and \emph{srij} (Guangyun, in Baxter's transcription), which glosses the first 'the law' and the second 'a model'.
\end{cousinweb}

\begin{grammarlens}[title={What you've built}]
One more word for saying what I do.
\end{grammarlens}

\subsection*{Your turn}
\begin{itemize}
\raggedright
  \item \emph{Say it:} \emph{lǜshī}
  \item \emph{Say it:} \emph{lǜshī}, once more
  \item \emph{Say it:} \emph{wǒ jiějie shì lǜshī} — My sister is a lawyer.
  \item \emph{Recall:} say the Mandarin for an occupation, then the Mandarin for an author, then say \emph{lǜshī} again
\end{itemize}

\begin{tcolorbox}[breakable,colback=teal!4,colframe=teal!35!black,title={Before you move on}]
What does lǜshī mean? (\textquotedblleft{}Lawyer\textquotedblright{}.) Say it once more.
\end{tcolorbox}
\section[jìzhě]{jìzhě — journalist; reporter}

\label{lesson:ZH-C259-jizhe}



\begin{quote}
Before the new one: say the Mandarin for an author, then the Mandarin for a lawyer.
\end{quote}

\subsection*{You'll want to know: jìzhě}
\emph{jìzhě} — \textquotedblleft{}journalist; reporter\textquotedblright{}.

\begin{quote}
\emph{tā shì yí ge jìzhě} — He's a journalist.
\end{quote}

Its characters come later; for now it is a word for the ear.

\begin{cousinweb}
\emph{jìzhě} is built from \emph{jì} 'to record' and \emph{zhě} 'one who', both inherited from Middle Chinese, \emph{kiH} and \emph{tsyaeX} (Guangyun, in Baxter's transcription), which glosses the first 'to note down' and calls the second a particle.
\end{cousinweb}

\begin{grammarlens}[title={What you've built}]
One more word for saying what I do.
\end{grammarlens}

\subsection*{Your turn}
\begin{itemize}
\raggedright
  \item \emph{Say it:} \emph{jìzhě}
  \item \emph{Say it:} \emph{jìzhě}, once more
  \item \emph{Say it:} \emph{tā shì yí ge jìzhě} — He's a journalist.
  \item \emph{Recall:} say the Mandarin for an author, then the Mandarin for a lawyer, then say \emph{jìzhě} again
\end{itemize}

\begin{tcolorbox}[breakable,colback=teal!4,colframe=teal!35!black,title={Before you move on}]
What does jìzhě mean? (\textquotedblleft{}Journalist; reporter\textquotedblright{}.) Say it once more.
\end{tcolorbox}
\section[\texorpdfstring{\zh{电工} (diàngōng)}{电工 (diàngōng)}]{\zh{电工} (diàngōng) — electrician}

\label{lesson:ZH-C259-diangong}



\begin{quote}
Before the new one: say the Mandarin for a lawyer, then the Mandarin for a journalist.
\end{quote}

\subsection*{You'll want to know: \zh{电工}}
\textbf{\zh{电工}} — \emph{diàngōng} — \textquotedblleft{}electrician\textquotedblright{}.

\begin{quote}
\textbf{\zh{他是电工}} \emph{tā shì diàngōng} — He's an electrician.
\end{quote}

\begin{cousinweb}
\textbf{\zh{电工}} (\emph{diàngōng}) is built from \zh{电} 'electricity' and \zh{工} 'a worker', both inherited from Middle Chinese, \emph{denH} and \emph{kuwng} (Guangyun, in Baxter's transcription), which glosses the first 'lightning'.
\end{cousinweb}

\begin{grammarlens}[title={What you've built}]
One more word for saying what I do.
\end{grammarlens}

\subsection*{Your turn}
\begin{itemize}
\raggedright
  \item \emph{Say it:} \emph{diàngōng}
  \item \emph{Say it:} \emph{diàngōng}, once more
  \item \emph{Say it:} \emph{tā shì diàngōng} — He's an electrician.
  \item \emph{Recall:} say the Mandarin for a lawyer, then the Mandarin for a journalist, then say \emph{diàngōng} again
\end{itemize}

\begin{tcolorbox}[breakable,colback=teal!4,colframe=teal!35!black,title={Before you move on}]
What does \zh{电工} mean? (\textquotedblleft{}Electrician\textquotedblright{}.) Say it once more.
\end{tcolorbox}
\section[(dialogue)]{Looking back: from Born and Married to What Do You Do?}

\label{lesson:ZH-R259-review-254-259}



\begin{quote}
Say the Mandarin for to be born and to get married.
\end{quote}

\subsection*{Your turn}
\begin{itemize}
\raggedright
  \item \emph{Say it:} to be born, to get married, to graduate, a wedding, daily life
  \item \emph{Say it:} to work part-time, wages, a colleague, to work overtime, the boss
  \item \emph{Say it:} to hold a view, a point of view, it looks as if, it may be that, different
  \item \emph{Say it:} a musical instrument, to play the drums, to perform on stage, a singer, to applaud
  \item \emph{Say it:} friendly, eager to help, patient with people, witty, personality
  \item \emph{Say it:} an occupation, an author, a lawyer, a journalist, an electrician
  \item \emph{Read it:} \zh{工} \zh{热} — each one aloud, then a word that holds it
\end{itemize}

\begin{tcolorbox}[breakable,colback=teal!4,colframe=teal!35!black,title={Before you move on}]
To be born? (\textbf{\zh{出生}}, \emph{chūshēng}.) To get married? (\emph{jiéhūn}.) To graduate? (\emph{bìyè}.) A wedding? (\emph{hūnlǐ}.) Daily life? (\emph{shēnghuó}.) To work part-time? (\textbf{\zh{打工}}, \emph{dǎgōng}.) Wages? (\emph{gōngzī}.) A colleague? (\textbf{\zh{同事}}, \emph{tóngshì}.) To work overtime? (\emph{jiābān}.) The boss? (\emph{lǎobǎn}.) To hold a view? (\emph{rènwéi}.) A point of view? (\emph{kànfǎ}.) It looks as if? (\textbf{\zh{看来}}, \emph{kànlai}.) It may be that? (\emph{yěxǔ}.) Different? (\textbf{\zh{不同}}, \emph{bùtóng}.) A musical instrument? (\emph{yuèqì}.) To play the drums? (\emph{dǎgǔ}.) To perform on stage? (\emph{yǎnchū}.) A singer? (\emph{gēshǒu}.) To applaud? (\emph{gǔzhǎng}.) Friendly? (\emph{yǒuhǎo}.) Eager to help? (\textbf{\zh{热心}}, \emph{rèxīn}.) Patient with people? (\emph{nàixīn}.) Witty? (\emph{yōumò}.) Personality? (\emph{xìnggé}.) An occupation? (\emph{zhíyè}.) An author? (\textbf{\zh{作家}}, \emph{zuòjiā}.) A lawyer? (\emph{lǜshī}.) A journalist? (\emph{jìzhě}.) An electrician? (\textbf{\zh{电工}}, \emph{diàngōng}.)
\end{tcolorbox}
